\section{El cálculo numérico}
\label{sec:numerico}

El cálculo numérico evalúa la ecuación tensorial adoptada sobre fondos concretos. El residuo $X$ controla la consistencia interna de cantidades construidas con ecuaciones relacionadas; su anulación no prueba de forma independiente la teoría ni su acción cuadrática.
\subsection{Estrategia y unidades}

El cálculo tiene tres etapas. Primero se resuelve el fondo cosmológico (la evolución de $\phi(t)$, $H(t)$ y $a(t)$) por separado para RG y para UG con difusión, con las mismas condiciones iniciales. Segundo, sobre cada fondo, se integra numéricamente la ecuación de cada modo tensorial $k$ desde dentro del horizonte, con condiciones de Bunch--Davies, hasta que el modo queda congelado fuera del horizonte. Tercero, se calcula $P_T(k)$ con la definición \eqref{eq:PTdef} y se guarda.

Se usan unidades $\MP=1$ (o sea $\kap=1$) y $m=1$, con el tiempo en unidades de $1/m$. Es exacto por homogeneidad: las ecuaciones \eqref{eq:H2UG}, \eqref{eq:KGQ}, \eqref{eq:modok} y \eqref{eq:S2v} no cambian de forma bajo $t\to t/m$, $H\to mH$, $V\to m^2V$, $Q\to m^2Q$, $k\to mk$. La masa física del inflatón, $m_{\rm phys}=6\times10^{-6}\MP$, entra solo al final, multiplicando $P_T$ por $m_{\rm phys}^2$ (porque $P_T\propto H^2/\MP^2$). Los números de $k$ que aparecen en las figuras están en unidades de $m$ con el factor de escala inicial $a_i=1$.

\subsection{El fondo en tiempo cósmico}

La variable independiente es $t$, y el estado es $(\phi,\dot\phi,N)$ con $N=\ln(a/a_i)$ y $\dot N=H$. Para RG el sistema es \eqref{eq:FriedPhi} y \eqref{eq:KG}, con $H=\sqrt{\rho/3}$ obtenido algebraicamente de Friedmann. Para UG es \eqref{eq:H2UG} para $H$ y \eqref{eq:KGQ} para $\ddot\phi$, con $Q(N)=Q_i\ee^{-\gamma N}$, de modo que $\dot Q=-\gamma QH$ (que es lo que entra en \eqref{eq:KGQ}).

Hay una decisión de diseño importante: el código \emph{no usa} la relación $\dot H=-\dot\phi^2/2\MP^2$ de \eqref{eq:HdotUG}. La cantidad $\dot H$, que hace falta para $\epsilon_1$ y para el término de masa que se describe abajo, se calcula derivando en el tiempo la ecuación de Friedmann de cada modelo y usando su ecuación del inflatón. Para UG, derivando \eqref{eq:H2UG},
\begin{equation}
6\MP^2H\dot H=\dot\phi\,(\ddot\phi+V')+\dot Q,
\label{eq:HdotCode}
\end{equation}
y análogamente en RG sin $\dot Q$. Así, \eqref{eq:HdotUG} queda como una \emph{consecuencia} que se puede contrastar con la solución integrada (control C3.2 de la sección~\ref{sec:resultados}) y no como un dato de entrada. Para terminar la inflación se usa el evento $\epsilon_1=1$, y las funciones $H(N)$, $\epsilon_1(N)$, $\phi(N)$ y el término de masa se guardan como splines cúbicos en $N$ sobre una malla de 20\,001 puntos.

\subsection{La ecuación de modos con el término de masa}
\label{sec:ecmodosnum}

Como control interno de implementación, el código integra una ecuación de modos que incluye explícitamente el término de masa. Se define
\begin{equation}
X\equiv-3H^2-2\dot H-\kap p+\Lb,
\label{eq:Xdef}
\end{equation}
que es el corchete de \eqref{eq:cancel}, con $\Lb=\Lambda$ en RG y $\Lb=\Lambda_0+Q(t)$ en UG. Si se mantiene $X$ en la acción, $S=\frac{1}{8\kap}\int[a^3\dot h_{ij}^2-a(\partial h_{ij})^2+2a^3X\,h_{ij}^2]$, la ecuación de movimiento es \etq{P}
\begin{equation}
\ddot h_k+3H\dot h_k+\Big(\frac{k^2}{a^2}-2X\Big)h_k=0,
\label{eq:modoX}
\end{equation}
que se reduce a \eqref{eq:modok} cuando $X=0$. La derivación analítica dice que $X=0$ por la ecuación de fondo; el código calcula $X$ con las cantidades de fondo de cada modelo y lo \emph{incluye} en la ecuación, como diagnóstico algebraico interno. C3.2 contrasta derivadas de splines; S1 contrasta la proyección TT lineal. Ninguno convierte este residuo algebraico en una prueba independiente de la acción.

Se integra en la variable $N$ y para $v=ah$ (o, más precisamente, para $\tilde v\equiv\sqrt{2k}\,v$, que es de orden uno al comienzo). En tiempo conforme, \eqref{eq:modoX} equivale a $v''+(k^2-a''/a-2Xa^2)v=0$. Con $\dd/\dd\eta=aH\,\dd/\dd N$, $a''/a=a^2H^2(2-\epsilon_1)$ y $\dd(aH)/\dd N=aH(1-\epsilon_1)$, y llamando $w\equiv\dd v/\dd N$, resulta \etq{P}
\begin{equation}
\frac{\dd w}{\dd N}+(1-\epsilon_1)\,w+\Big[\frac{k^2}{(aH)^2}-(2-\epsilon_1)-\frac{2X}{H^2}\Big]v=0.
\label{eq:modoN}
\end{equation}
El cálculo de $a''/a$ es directo: $a'=a^2H$, luego $a''=a\,\tfrac{\dd}{\dd t}(a^2H)=a^3(2H^2+\dot H)$, es decir $a''/a=a^2(2H^2+\dot H)=a^2H^2(2-\epsilon_1)$. Esta forma es cómoda porque la ecuación solo depende del fondo a través de $\epsilon_1(N)$, $H(N)$ y $X(N)/H^2$.

\subsection{Condiciones iniciales y el espectro}

Cada modo se arranca en el instante $N_s$ en que $k/(aH)=100$ (dentro del horizonte), con la condición de Bunch--Davies en su versión adiabática (WKB),
\begin{equation}
v=\frac{1}{\sqrt{2\omega}},\qquad v'=-i\,\omega\,v,\qquad \omega^2=k^2-\frac{a''}{a}-2Xa^2,
\label{eq:BDwkb}
\end{equation}
que se reduce a $v=\ee^{-ik\eta}/\sqrt{2k}$ para $k\to\infty$ \etq{S}. Se desprecia la derivada de $\omega$; el error relativo de esa aproximación es del orden de $(aH/k)^3\sim10^{-6}$ para $k/(aH)=100$, y se verifica numéricamente en la sección~\ref{sec:resultados}. Se integra hasta $k/(aH)=10^{-3}$ (o hasta el fin de la inflación si ocurre antes), donde el modo está congelado con un error $\sim(k/aH)^2=10^{-6}$. El espectro se calcula con \eqref{eq:PTdef},
\begin{equation}
P_T=\frac{4k^3|v|^2}{\pi^2a^2}=\frac{2k^2|\tilde v|^2}{\pi^2a^2}\qquad(\MP=1).
\label{eq:PTcode}
\end{equation}
La integración usa el método de Dormand--Prince de orden 8 (\texttt{DOP853}) de \texttt{scipy} con tolerancia relativa $10^{-10}$, y los modos se calculan en paralelo.

\subsection{Decisiones de modelado}
\label{sec:decisiones}

Las decisiones que fijan el cálculo se resumen en la tabla~\ref{tab:decisiones}. La mayoría se tomaron con el usuario del trabajo al empezar el cálculo; las que no, se marcan como valores por defecto. Interesa detenerse en dos.

La primera es la elección de $\Lambda_0$. Cuando se comparan dos modelos con las mismas condiciones iniciales, hay que decidir qué significa ``las mismas'' (sección~\ref{sec:fondo}). Se eligió fijar el mismo $H_{\rm ini}$ en RG y en UG, con el mismo $(\phi_i,\dot\phi_i)$. Por \eqref{eq:Lambda0}, eso obliga a $\Lambda_0=-Q_i/\MP^2$, y se vuelve negativa. Como se verá, esta decisión tiene consecuencias grandes sobre la magnitud del efecto (sección~\ref{sec:fondores}). La segunda es la duración de la inflación. Se fijó que UG dure $N_f=100$ e-folds, valor que aparece en la literatura de UG con difusión: \cite{piccirilli23} usa $N_f=100$ en su primer escenario y afirma que $N_f\ge65$ es lo habitual, y \cite{leon22} ilustra con $N_f=70$ y supone al menos $\sim102$ e-folds en una condición de su análisis (lecturas parciales de esos textos). Como consecuencia, $\phi_i$ se determina por una búsqueda de raíz para que UG dure exactamente $100$ e-folds, y RG, con el mismo $\phi_i$, dura más.

\begin{table}[t]
\centering\small
\renewcommand{\arraystretch}{1.25}
\begin{tabular}{@{}p{4.0cm}p{5.2cm}p{4.2cm}@{}}
\toprule
\textbf{Decisión} & \textbf{Valor} & \textbf{Origen}\\
\midrule
Acoplamiento de $Q$ & solo al inflatón, \eqref{eq:KGQ} & decisión del usuario; coherente con un campo solo si $Q=Q(\phi)$ (sección~\ref{sec:acoplamiento})\\
Forma de $Q$ & $Q(N)=Q_i\ee^{-\gamma N}$, $Q_i/V(\phi_i)=0.1$, $\gamma=0.1$ & decisión del usuario\\
Potencial y masa & $V=\tfrac12m^2\phi^2$, $m=6\times10^{-6}\MP$; y Starobinsky \eqref{eq:Vstaro} con $M=1.14\times10^{-5}\MP$ (sección~\ref{sec:potenciales}) & decisión del usuario; $M$ calibrada con $A_s$\\
Definición de $P_T$ & $\langle h_{ij}h_{ij}\rangle$, $e\cdot e=2$ & decisión del usuario\\
Condiciones iniciales & mismos $\phi_i$, $\dot\phi_i$, $H_i$; $\Lambda_0$ por \eqref{eq:Lambda0} & decisión del usuario\\
Duración de UG & $N_f=100$ e-folds $\Rightarrow\phi_i=25.386\,\MP$ (cuadrático), $7.660\,\MP$ (Starobinsky) & criterio: literatura de UG\\
Criterio de comparación & a igual $k$ (sección~\ref{sec:resultados}) y a igual $N_*$ (sección~\ref{sec:potenciales}, principal) & decisión del usuario\\
Velocidad inicial & slow-roll de RG: $\dot\phi_i=-V'/\sqrt{3V}$ & por defecto\\
Arranque y medición del modo & $k/aH=100$ y $10^{-3}$ & por defecto\\
Malla de $k$ & 40 modos equiespaciados en $\ln k$; mismo $k$ comóvil en RG y UG, $a_i=1$ & por defecto\\
Tolerancias & modos $10^{-10}$; fondo $10^{-12}$ & por defecto\\
\bottomrule
\end{tabular}
\caption{Decisiones de modelado y de cálculo numérico.}
\label{tab:decisiones}
\end{table}

\subsection{Estructura del código y reproducibilidad}

El código está organizado en módulos pequeños: \texttt{config.py} (parámetros), \texttt{background.py} (fondo de RG y de UG, el término de masa), \texttt{modes.py} (ecuación de modos, condiciones iniciales y $P_T$), \texttt{run\_spectrum.py} (orquestación y salida), \texttt{checks.py} (los controles de la sección siguiente) y \texttt{make\_figures.py} (las figuras a partir de los resultados guardados). Cada función lleva en su documentación la ecuación que implementa, y el Apéndice~\ref{ap:codigo} lista la correspondencia completa. El cálculo es determinista, corre en unos pocos segundos, y hay un notebook (\texttt{explorador\_PT.ipynb}) que permite ejecutarlo paso a paso.
